\documentclass[a4paper,12pt]{article}

\usepackage[utf8]{inputenc}
\usepackage{amsmath}
\usepackage{graphicx}
\usepackage{booktabs}
\usepackage{listings}
\usepackage{xcolor}
\usepackage{geometry}

\geometry{margin=1in}

\title{Labwork 5: Neural Network with Backpropagation}
\author{Your Name}
\date{\today}

\definecolor{codegreen}{rgb}{0,0.6,0}
\definecolor{codegray}{rgb}{0.5,0.5,0.5}

\lstset{
    basicstyle=\ttfamily\small,
    breaklines=true,
    frame=single,
    numbers=left,
    numberstyle=\tiny\color{codegray},
    keywordstyle=\color{blue},
    commentstyle=\color{codegreen},
    showstringspaces=false
}

\begin{document}

\maketitle

\tableofcontents

\newpage

% ==================================================
\section{Introduction}

Artificial Neural Networks (ANNs) are computational models inspired by the human brain. 
They are widely used in machine learning for solving classification and regression problems.

In this labwork, a neural network was implemented completely from scratch using Python, without using machine learning libraries such as TensorFlow or PyTorch.

The objective of this labwork is:
\begin{itemize}
    \item Implement feedforward propagation
    \item Implement backpropagation
    \item Implement gradient descent optimization
    \item Train the network using datasets stored in CSV files
    \item Evaluate the neural network on different datasets
\end{itemize}

The experiments were performed on:
\begin{itemize}
    \item XOR dataset
    \item Loan approval dataset
    \item House price-size dataset
\end{itemize}

% ==================================================
\section{Neural Network Architecture}

The neural network is composed of:
\begin{itemize}
    \item Input layer
    \item Hidden layer
    \item Output layer
\end{itemize}

Each neuron contains:
\begin{itemize}
    \item Weights
    \item Bias
    \item Activation function
\end{itemize}

The sigmoid activation function was used:

\[
\sigma(x) = \frac{1}{1 + e^{-x}}
\]

The derivative of sigmoid is:

\[
\sigma'(x) = \sigma(x)(1 - \sigma(x))
\]

% ==================================================
\section{Feedforward Propagation}

During feedforward propagation, the inputs are multiplied by weights and added with biases.

The weighted sum is computed as:

\[
z = \sum_{i=1}^{n} x_i w_i + b
\]

Then the sigmoid activation function is applied:

\[
a = \sigma(z)
\]

The outputs from one layer become the inputs of the next layer.

% ==================================================
\section{Backpropagation}

Backpropagation is used to compute gradients and update the weights of the network.

The error of the output layer is:

\[
Error = Expected - Output
\]

The delta value for output neurons is:

\[
\delta = (Expected - Output)\sigma'(Output)
\]

For hidden layers:

\[
\delta_j = \sigma'(o_j)\sum_k w_{jk}\delta_k
\]

Weights are updated using gradient descent:

\[
w = w + \eta \delta x
\]

Biases are updated as:

\[
b = b + \eta \delta
\]

where:
\begin{itemize}
    \item $\eta$ is the learning rate
    \item $\delta$ is the neuron delta
    \item $x$ is the input value
\end{itemize}

% ==================================================
\section{Implementation}

The neural network was implemented in Python using the following classes:

\begin{itemize}
    \item \texttt{Neuron}
    \item \texttt{Layer}
    \item \texttt{NeuralNetwork}
\end{itemize}

The implementation supports:
\begin{itemize}
    \item Feedforward propagation
    \item Backpropagation
    \item Gradient descent training
    \item Weight saving/loading
\end{itemize}

The following Python libraries were used:
\begin{itemize}
    \item math
    \item random
    \item csv
\end{itemize}

No deep learning framework was used.

% ==================================================
\section{XOR Dataset Experiment}

The XOR dataset is a classical non-linear classification problem.

\subsection{Dataset}

\begin{center}
\begin{tabular}{ccc}
\toprule
$x_1$ & $x_2$ & Output \\
\midrule
0 & 0 & 0 \\
0 & 1 & 1 \\
1 & 0 & 1 \\
1 & 1 & 0 \\
\bottomrule
\end{tabular}
\end{center}

\subsection{Training Parameters}

\begin{itemize}
    \item Learning rate: 0.5
    \item Epochs: 10000
    \item Architecture: 2-2-1
\end{itemize}

\subsection{Training Result}

The training error decreased significantly during training:

\begin{verbatim}
Epoch 0 Error = 1.511978
Epoch 9000 Error = 0.001426
\end{verbatim}

Final predictions:

\begin{center}
\begin{tabular}{ccc}
\toprule
Input & Expected & Predicted \\
\midrule
[0,0] & 0 & 0.016163 \\
[0,1] & 1 & 0.981511 \\
[1,0] & 1 & 0.981505 \\
[1,1] & 0 & 0.017538 \\
\bottomrule
\end{tabular}
\end{center}

The neural network successfully learned the XOR function.

% ==================================================
\section{Loan Dataset Experiment}

\subsection{Dataset Description}

The loan dataset contains:
\begin{itemize}
    \item Experience
    \item Salary
    \item Loan approval result
\end{itemize}

The objective is to predict whether a loan should be approved.

\subsection{Data Normalization}

Input values were normalized to improve training performance.

\subsection{Training Parameters}

\begin{itemize}
    \item Learning rate: 0.5
    \item Epochs: 10000
    \item Architecture: 2-2-1
\end{itemize}

\subsection{Results}

The error decreased during training:

\begin{verbatim}
Epoch 0 Error = 5.098626
Epoch 9000 Error = 1.003566
\end{verbatim}

Sample predictions:

\begin{center}
\begin{tabular}{cccc}
\toprule
Experience & Salary & Expected & Predicted \\
\midrule
1.0 & 0.4 & 1 & 0.999965 \\
0.333 & 0.4 & 0 & 0.000001 \\
0.666 & 0.8 & 1 & 0.999996 \\
0.333 & 0.8 & 0 & 0.999823 \\
\bottomrule
\end{tabular}
\end{center}

The network achieved good classification performance, although some predictions were imperfect due to the small dataset size.

% ==================================================
\section{House Price Dataset Experiment}

\subsection{Dataset Description}

The dataset contains:
\begin{itemize}
    \item House size
    \item House price
\end{itemize}

The objective is to predict house price from house size.

\subsection{Normalization}

Prices were normalized between 0 and 1 because the sigmoid activation function outputs values in this range.

\subsection{Training Parameters}

\begin{itemize}
    \item Learning rate: 0.5
    \item Epochs: 10000
\end{itemize}

\subsection{Results}

Training error:

\begin{verbatim}
Epoch 0 Error = 1.400398
Epoch 9000 Error = 0.009570
\end{verbatim}

Predictions:

\begin{center}
\begin{tabular}{ccc}
\toprule
Size & Real Price & Predicted Price \\
\midrule
10 & 55 & 60.36 \\
20 & 80 & 73.03 \\
40 & 100 & 100.96 \\
60 & 120 & 126.10 \\
80 & 150 & 140.78 \\
\bottomrule
\end{tabular}
\end{center}

The neural network successfully approximated the relationship between house size and price.

% ==================================================
\section{Conclusion}

In this labwork, a neural network was successfully implemented from scratch using Python.

The implementation included:
\begin{itemize}
    \item Feedforward propagation
    \item Backpropagation
    \item Gradient descent optimization
    \item CSV dataset processing
\end{itemize}

The neural network successfully learned:
\begin{itemize}
    \item XOR classification
    \item Loan approval classification
    \item House price prediction
\end{itemize}

The experiments demonstrated that backpropagation can effectively optimize weights and reduce prediction error.

\end{document}